\section{Technical Decisions and Dependencies}
\label{sec:technical-decisions}

The \tradesq{} Mobile Application uses Expo, React Native, and TypeScript.
This combination provides a shared codebase for Android and iOS while retaining access to device capabilities through Expo and React Native platform APIs.
Flutter and Ionic were considered as alternative cross-platform approaches.
React Native was selected because the project team has experience with TypeScript and React-based development, and because the Expo ecosystem provides suitable integrations for routing, local storage, secure storage, image selection, localisation, and native application builds.

The application follows a layered Clean Architecture variant with presentation, application, domain, and infrastructure layers.
This approach was selected to separate screens and reusable user-interface components from application use cases, domain contracts, and technical integrations.
A simpler feature-only structure with direct API calls from screens was considered but rejected because it would couple interface code to the Back-end API, local persistence, and device services.
The layered structure supports isolated testing, replacement of repository implementations, and clear dependency direction.

\begin{htable}[
    caption = {Key third-party libraries and technical decisions},
    label = {tab:libraries},
]
    {
        colspec = {p{0.22\textwidth} p{0.24\textwidth} X},
    }
        \textbf{Technology} & \textbf{Purpose} & \textbf{Selection rationale and alternatives} \\
        Expo Router & File-based routing with Stack and tab navigation. & Expo Router was selected because it integrates with Expo and keeps route definitions close to route files. Manually configured React Navigation was considered but would require more navigation configuration. \\
        TanStack React Query & Server-state retrieval, mutations, cache management, and query persistence. & React Query was selected because it provides a consistent query and mutation model with cache invalidation and persisted query support. Manual \texttt{useEffect} requests, Context-only state, Zustand, and Redux Toolkit Query were considered. \\
        React Hook Form and Zod & Form state management and schema-based validation. & This combination reduces form boilerplate and provides typed validation schemas. Formik, Yup, and manual component state were considered as alternatives. \\
        Expo SecureStore & Encrypted storage of session tokens. & SecureStore was selected because it uses platform-backed protected storage. AsyncStorage and SQLite were rejected for tokens because they are intended for non-sensitive application data. \\
        Expo SQLite and Drizzle ORM & Local persistence of cached data and listing drafts. & SQLite provides structured persistent storage, while Drizzle provides typed schema and query definitions. AsyncStorage alone was considered but is less suitable for structured listing and draft data. \\
        React Native NetInfo & Network-connectivity status. & NetInfo provides explicit connectivity information for the offline banner and the prevention of unavailable actions. Relying only on failed API requests was considered but provides less immediate feedback. \\
        Expo Image Picker and React Native Share & Image selection and native sharing. & These libraries provide cross-platform access to device media selection and sharing capabilities. Custom native implementations were not selected because they would increase platform-specific maintenance. \\
        react-i18next & English and Dutch localisation. & The library centralises translations and language state. Manually maintained translation objects were considered but would make localisation logic harder to maintain. \\
        NativeWind & Utility-based styling for React Native components. & NativeWind was selected because it provides Tailwind-style utility classes and keeps component styling close to the relevant interface code. React Native \texttt{StyleSheet} objects and styled-components were considered as alternatives. \\
        react-native-reusables & Reusable accessible interface primitives. & The library provides composable low-level controls that can be styled with NativeWind. NativewindUI was considered as a prebuilt NativeWind-compatible alternative but was not selected because it is a paid dependency. The selected open-source libraries provide the required flexibility. \\
        Lucide React Native & Icon library for application controls and navigation. & Lucide was selected because it provides a consistent and broad icon set for React Native. Expo Vector Icons and custom SVG assets were considered as alternatives. \\
\end{htable}

